% Generated draft: requires mathematical and source review.
\documentclass[11pt]{article}
\usepackage[margin=25mm]{geometry}
\usepackage{fontspec}
\setmainfont{Latin Modern Roman}
\usepackage{amsmath,amssymb,mathtools,unicode-math}
\setmathfont{Latin Modern Math}
\usepackage{enumitem,needspace,xurl,hyperref}
\setlength{\parindent}{0pt}
\setlength{\parskip}{5pt}
\setlength{\emergencystretch}{3em}
\newcommand{\N}{\mathbb{N}}
\newcommand{\Z}{\mathbb{Z}}
\newcommand{\Q}{\mathbb{Q}}
\newcommand{\R}{\mathbb{R}}
\newcommand{\C}{\mathbb{C}}
\newcommand{\F}{\mathbb{F}}
\newcommand{\A}{\mathbb{A}}
\newcommand{\PP}{\mathbb{P}}
\newcommand{\E}{\mathbb{E}}
\newcommand{\eps}{\varepsilon}
\newcommand{\dd}{\,\mathrm{d}}
\begin{document}
\section*{UnsolvedMath ID: 30003826}
\textbf{Source problem code:} catalogue entry
\subsection{Definitions and mathematical statement}
The supplied record asks for analogues of three assertions for the enveloping algebras of the full Witt algebra and the Virasoro algebra. It does not specify the ground field, the precise meaning of a quadratic expression, or unambiguously which quotient is intended. The following gives a conditional formulation and makes these unresolved qualifications explicit.

For the conventional characteristic-zero interpretation, let \(k\) be a field with \(\operatorname{char}k=0\). Define the full Witt algebra by
\[
W=\bigoplus_{n\in\Z}ke_n,
\qquad [e_m,e_n]=(n-m)e_{m+n}.
\]
Define the Virasoro algebra by
\[
V=\left(\bigoplus_{n\in\Z}ke_n\right)\oplus kc,
\qquad [c,e_n]=0,
\]
\[
[e_m,e_n]=(n-m)e_{m+n}
 +\delta_{m+n,0}\frac{m^3-m}{12}c,
\]
where \(\delta_{a,0}\) is the Kronecker delta. These presentations specify a conventional interpretation; the supplied record does not establish that characteristic zero, this normalization, or any particular central specialization is intended. For \(\mathfrak g\in\{W,V\}\), write \(A=U_k(\mathfrak g)\) for its universal enveloping algebra.

For a nonzero unital \(k\)-algebra \(B\), define its Gelfand--Kirillov dimension by
\[
\operatorname{GKdim}_k B
 =\sup_E\limsup_{n\to\infty}
 \frac{\log\bigl(\dim_k E^n\bigr)}{\log n},
\]
where \(E\) ranges over finite-dimensional \(k\)-linear subspaces of \(B\) containing \(1\), and \(E^n\) is the linear span of products of \(n\) elements of \(E\). Adopt \(\operatorname{GKdim}_k(0)=-\infty\), so that the zero quotient satisfies a finite-growth requirement.

Under the interpretation that ``an ideal has a quotient'' means the quotient of \(A\) by that ideal, the three questions are whether, for each \(\mathfrak g\in\{W,V\}\), the following assertions hold:
\begin{enumerate}
\item For every nonzero two-sided ideal \(I\) of \(A\),
\[
\operatorname{GKdim}_k(A/I)<\infty.
\]
\item Every ascending sequence of two-sided ideals
\[
I_1\subseteq I_2\subseteq I_3\subseteq\cdots
\]
stabilizes: there exists \(N\geq1\) such that \(I_n=I_N\) for all \(n\geq N\).
\item If a two-sided ideal \(I\) contains a quadratic expression in the intended generators, then
\[
\operatorname{GKdim}_k(A/I)<\infty.
\]
\end{enumerate}
The third assertion still requires a specified class \(\mathcal Q\subseteq A\) of quadratic expressions; with that class fixed, its hypothesis is \(I\cap\mathcal Q\neq\varnothing\). For example, the standard enveloping-algebra filtration is
\[
F_dA=\operatorname{span}_k\{x_1\cdots x_r:0\leq r\leq d,\ x_i\in\mathfrak g\}.
\]
Interpreting quadratic as exact filtration degree two would give \(\mathcal Q=F_2A\setminus F_1A\). This is only a possible interpretation: the record does not determine whether lower-degree terms are allowed, whether homogeneity is required, which generators may occur, or whether a nonzero element is required. In particular, it does not specify the role of the central generator \(c\).

The original wording could instead refer to a quotient of the ideal itself. No such interpretation is adopted here without qualification. The report's referenced preceding three assertions are not reproduced with their hypotheses, so this conditional formulation requires comparison with the original report before it can serve as a definitive statement.
\subsection{Short English statement}
For the full Witt algebra and the Virasoro algebra, ask whether every nonzero two-sided ideal yields an enveloping-algebra quotient of finite Gelfand--Kirillov dimension, whether ascending chains of two-sided ideals stabilize, and whether containing a quadratic expression forces such a finite-dimensional-growth quotient. The ground field, the intended quotient, and the meaning of quadratic require clarification. The supplied open label is dated curation, not an independently verified conclusion.
\subsection{Sources}
\begin{itemize}
\item Supplied local UnsolvedMath registry record, permanent ID \(30003826\), titled \emph{Ideal Theory of Witt and Virasoro Enveloping Algebras}; Oberwolfach Reports Open Problems collection; exact source problem code \texttt{catalogue entry}. The record gives proposed year 2018 and names no individual proposer.
\item Supplied report citation: \emph{Interactions between Algebraic Geometry and Noncommutative Algebra} (2019), Oberwolfach Reports, pp.\ 1465--1515, DOI \texttt{10.4171/owr/2018/24}. Supplied source URL and literature-source link, labeled ``original OWR report'': \url{https://doi.org/10.4171/owr/2018/24}. The background describes Oberwolfach Reports as documenting mathematical workshops and their open-problem discussions. No exact problem page or individual attribution is supplied.
\item Literature assessment checked at \texttt{2026-08-21}: the supplied curation labels the problem open, describes the proposed analogues for the full Witt and Virasoro algebras as future work, and reports finding no proof in the primary literature checked. The listed checked source is the original OWR report linked above. This is supplied dated curation, not an independently verified literature assessment; the background explicitly requests independent verification of time-sensitive status and cited claims.
\item The registry labels the statement \texttt{corrected\_verified} and reports a manual semantic, typographic, and LaTeX review claiming that the clean statement preserves the target. These are supplied editorial assertions; this draft does not independently verify them or establish mathematical correctness.
\end{itemize}
\end{document}
